\newpage
\topic{8 continued}{A complete SIMPLE pass on the same two cells}
Use the closed two-cell geometry of Section 7: one shared horizontal velocity
$u$, cell spacing $h=0.5$, and zero normal velocity at every external wall.
There is only one unknown velocity, so the column $q$ has just the entry $u$.
The two pressures are $p_1$ and $p_2$.

For this algebraic exercise, prescribe the original, volume-normalized momentum
row
\[
7u=2-\frac{p_2-p_1}{0.5}.
\]
Thus the unrelaxed center coefficient is $a^0=7$ and the known source is
$b=2$. These specified numbers isolate the coupling operations; they are not
a spatially resolved cavity calculation. Start from $u^0=0$ and
$p_1^0=p_2^0=0$, and select $\alpha_u=0.7$, $\alpha_p=0.3$.

\textbf{1. Relax the momentum row.} At the first pass, $k=0$:
\[
a^{\rm rel}=7/0.7=10,\qquad
b^{\rm rel}=2+(1/0.7-1)7u^0=2.
\]
\textbf{2. Solve the predictor.} The current pressure gradient is zero, so
\[
u^*=\frac{2-(0-0)/0.5}{10}=0.2,\qquad d=1/10=0.1.
\]
The values $u^*=0.2$ and $d=0.1$ used in Section 7 now follow from the same
momentum row and current estimates.

\textbf{3. Calculate the cell imbalances.} Positive $u^*$ leaves cell 1 and
enters cell 2:
\[
Dq^*=\begin{bmatrix}u^*/h\\-u^*/h\end{bmatrix}
=\begin{bmatrix}0.4\\-0.4\end{bmatrix}.
\]
\textbf{4. Solve the pressure correction.} Use the matrix from Section 7:
\[
0.4\begin{bmatrix}1&-1\\-1&1\end{bmatrix}
\begin{bmatrix}p'_1\\p'_2\end{bmatrix}
=\begin{bmatrix}-0.4\\0.4\end{bmatrix}.
\]
Its first row gives $p'_1-p'_2=-1$. Impose the zero-mean reference
$p'_1+p'_2=0$, giving $p'_1=-0.5$, $p'_2=0.5$. The second mass-balance row
then follows as well; it is dependent because the cells share their only
interior face and all external normal flow rates are zero.
\newpage
\topic{8 continued}{Correct both fields, then check the original equations}
\textbf{5. Update velocity and stored pressure.} Apply the full face correction:
\[
u^1=u^*-d\frac{p'_2-p'_1}{h}
=0.2-0.1\frac{0.5-(-0.5)}{0.5}=0.
\]
Store the under-relaxed pressure for the next predictor:
\[
p_1^1=0+0.3(-0.5)=-0.15,\qquad
p_2^1=0+0.3(0.5)=0.15.
\]
The superscript 1 means one completed steady coupling pass, not a physical
time step. Both cell continuity residuals are now zero.

\textbf{6. Check momentum using the original row.} Its scalar residual is
\[
r_m=7u-b+\frac{p_2-p_1}{h}.
\]
For the updated fields this is
\[
r_m^1=7(0)-2+\frac{0.15-(-0.15)}{0.5}=-1.4.
\]
The original momentum equation is not yet satisfied with the stored pressure,
even though continuity is exact. This is why the next SIMPLE pass is needed.

At pass 2, use $u^1=0$ and $p^1=(-0.15,0.15)$ as inputs. The relaxed source
remains 2, and
\[
u^*=\frac{2-0.30/0.5}{10}=0.14.
\]
The next pressure correction is $p'=(-0.35,0.35)$. Correcting velocity fully
again gives zero; storing $0.3p'$ gives $p^2=(-0.255,0.255)$ and
$r_m^2=-0.98$. Repeating the same operations gives:
\begin{center}\begin{tabular}{rrrrr}\toprule
Completed pass & $u^*$ & Corrected $u$ & Stored $(p_1,p_2)$ & $|r_m|$\\\midrule
1 & 0.200 & 0 & $(-0.150,0.150)$ & 1.400\\
2 & 0.140 & 0 & $(-0.255,0.255)$ & 0.980\\
3 & 0.098 & 0 & $(-0.3285,0.3285)$ & 0.686\\\bottomrule
\end{tabular}\end{center}
The exact constrained steady solution is $u=0$, $p_1=-0.5$, $p_2=0.5$.
This one-velocity example has no off-diagonal velocity coupling, so its
diagonal pressure-response approximation is exact for a frozen row.
The chosen pressure under-relaxation makes the stored pressures approach
the solution over successive passes. Real cavity rows also contain neighbor
coupling and velocity-dependent convection coefficients.
